\chapter{118}



«Dieses kleine Biest.»

Louis murmelte die Worte vor sich hin und ging ununterbrochen in seiner
Wohnung umher. Der kurze Auftritt des Mädchens hatte ihn komplett
durcheinander gewirbelt. Da hatte ihm auch \emph{Ed Sheerans} positive
Energie nichts genützt.

Noch lag Valentinas Brief verschlossen auf der Ablage im Flur. Er hatte
ihn mehrmals hochgehoben, aber immer gleich wieder zurückgelegt, als
wäre er vergiftet. Irgend etwas hinderte ihn daran, ihn zu öffnen. Kam
er sich nur abserviert vor oder war er einfach unfähig, sich
Veränderungen anzupassen? Er war definitiv zu einem Nervenbündel
geworden.
Joh würde erst gegen Abend zurück sein.
Louis musste in Bewegung kommen. Den Brief fasste er nicht mehr an.

Ohne weitere Gedanken holte er den Staubsauger aus dem Schrank, steckte
ihn ein und kurvte wie ein Versessener durch die Räume, als gelte es,
einen Preis für Reinlichkeit zu gewinnen. Dann bezog er sein Bett
frisch, pumpte die Schlafmatratze praller, stülpte in hohem Tempo
frische Bettbezüge über Ersatzdecke und Kissen und schrubbte an Dusche
und Lavabo herum. Erst in der Küche, als ihm die Milchtüte aus dem
Kühlschrank entgegen sprang und zu Boden knallte und sich die Flüssigkeit
verteilte, kam er wieder zu sich. Entgeistert schaute er zu, wie sich
verschieden breite Milchflüsse fortbewegten, sich den Plattenfugen
entlang verteilten und nach und nach stoppten.
Er setzte sich an den Tisch. 
Er war unverbesserlich. Wie hatte ihm die Begegnung mit diesem Teenager
nur so zusetzen können. Was war denn schon passiert? Und, verdammt, wie
war das nun mit seiner Flexibilität bei unverhofften Wendungen?
Louis fuhr sich über die Stirn. Dann musste er lachen. Seine Aktion
hatte einzig und alleine dazu gedient, seine Hilflosigkeit abzuschalten.
Sein Lachen wurde lauter. 

Bevor er den Boden wischte, machte er sich einen Kaffee.

\sectionbreak

Mit Johs Rückkehr war auch das Nachtessen bereit. Zumindest das Kochen hatte Louis
Spass gemacht.

«Das riecht wie in einem 5-Sterne-Restaurant.»

Johs erste Worte beim Eintreten in die Wohnung hörten sich erfrischend
an.

«Ja, ja,» rief Louis ausgelassen, «Ossobuco, meine Königsdisziplin.»

Noch vor dem Essen warfen sie eine Münze. Es war Johs
Idee.

«Also, bei Kopf erzählst du zuerst von deiner Begegnung mit Valentina, bei Zahl fange ich an und erstatte dir Rapport über meine Mailand-Tour mit Josephine.»

Joh gewann. Maman musste alle Register gezogen haben. Während sie 
assen, erzählte er mit Begeisterung, wohin sie ihn geführt hatte. 
Von ihren vielen Informationen zu den einzelnen Orten. Erzählte mit Ehrfurcht von ihren
Gesprächen und nannte seine Mutter eine wunderbare Frau. Louis hätte ihm
noch lange zuhören mögen.

«Schön, Joh, wie du sie erlebt hast. Ja, sie ist ein toller Mensch, eine 
Mutter, wie man sie sich nur wünschen kann. Ich liebe sie.»

Joh sah ergriffen aus. Für einen Augenblick war es still.
Dann klopfte er sich auf die Brust und schloss die Augen.

«Dein Ossobuco ist unerreicht, Louis, wow, diese Zubereitung.»

«Danke, freue mich, dass es dir schmeckt, ja, es ist gelungen. Ich bin zufrieden.»

Schliesslich griff Joh sein Glas, lehnte sich zurück und setzte seinen wachen Blick auf.

«Und nun, Louis, bist du dran.»

Louis fasste mit wenigen Worten die Bar-Szene zusammen. Erzählte von
seiner Enttäuschung, seiner Wut auch, schliesslich, dass nun wenigstens
die Wohnung sauber sei und blieb bei Valentinas Brief stecken.
Joh wollte gerade die Gabel zum Mund führen, als er inne hielt. 

«Wow, und du hast den Brief noch nicht gelesen, wirklich?»

Louis nickte und zog seine Schultern hoch.
Dann stand er auf, ging in den Flur und setzte sich mit dem Umschlag in
der Hand wieder an den Tisch.

«La voilà.»

Louis legte den Brief zwischen Weinflasche und Salatschüssel.
Für einen Moment schauten sie ihn beide an.

«Spannend,» Joh suchte nach Worten, was nicht oft vorkam, «und nun?»

«Ich weiss nicht.»

«Wovor fürchtest du dich?»

«Vielleicht, dass ich mich noch schuldiger fühle als eh schon?»

«Hm.»

Joh fuhr sich über die Oberlippe.

«Was könnte als beste Variante drin stehen?»

Louis musste nicht lange nachdenken.

«Dass sie wirklich krank ist, mir aber unbedingt mitteilen will, dass
sie mir vergibt. So etwas.»

Louis wollte nach dem Brief greifen, zog aber seine Hand gleich wieder
zurück.

«Ich muss ihn lesen, Joh. Ich muss mich frei machen. Morgen erwartet
mich etwas viel Grösseres. Ich habe einen Riesenbammel vor der Begegnung
mit Giorgia.»

Beim Aussprechen des Namens fuhr Louis zusammen.

«Gut, wie kann ich dir helfen, Louis?»

Louis starrte auf die Tischmitte. Die Sekunden vergingen. Joh verschränkte die
Arme.
Mit einem Mal hob Louis den Kopf und blickte Joh eindringlich an.

«Bitte, lies \emph{du} ihn. Und, und fass mir den Inhalt in drei Worten
zusammen, ok?»

Joh nahm den Brief an sich und schlitzte ihn mit dem Buttermesser auf.

«Mach ich gern für dich.»

Louis beobachtete, wie Johs Augen über die Zeilen flogen. Sein
Ausdruck verriet vorerst nichts über den Inhalt. Dann
glaubte Louis, ein feines Lächeln auf seinen Lippen zu sehen.
Joh faltete den Brief und steckte ihn zurück in den Umschlag. Er fasste
über den Tisch nach Louis’ Hand.

«Ok, in drei Worten, Louis: \emph{Valentina ist verliebt}.»

Sie assen schweigend weiter.
